% -*- coding: utf-8 -*-
% !TEX program = xelatex
\documentclass[plain,noanswer,medmath]{../../template/jnuexam}
\usepackage{../../template/kysx}

\begin{document}


\examtitle{2003}{4}


\loadproblems{3}{2003P3}%载入试卷三


\makepart{填空题}{本题共6小题，每小题4分，满分24分}


\begin{problem}
极限$\lim_{x \to 0} [1 + \ln(1 + x)]^{\frac{2}{x}} =$ \fillin{}．
\end{problem}


\begin{problem}
$\int_{-1}^1 (|x| + x) \e^{- |x|} \dx =$ \fillin{}．
\end{problem}


\useproblem{3}{1}{3}%【同试卷三第一(3)题】


\begin{problem}
设$A$，$B$均为三阶矩阵，$E$是三阶单位矩阵．已知$AB = 2A + B$，$B = \begin{pmatrix}
  2 & 0 & 2 \\
  0 & 4 & 0 \\
  2 & 0 & 2
\end{pmatrix}$，则$(A - E)^{-1} =$ \fillin{}．
\end{problem}


\useproblem{3}{1}{4}%【同试卷三第一(4)题】


\begin{problem}
设随机变量$X$和$Y$的相关系数为$0.5$，$EX = EY = 0$，$EX^2 = EY^2 = 2$，则$E(X + Y)^2 =$ \fillin{}．
\end{problem}


\makepart{选择题}{本题共6小题，每小题4分，满分24分}


\begin{problem}
曲线$y = x\e^{\frac{1}{x^2}}$\pickout{}
\begin{abcd}
\item 仅有水平渐近线．
\item 仅有铅直渐近线．
\item 既有铅直又有水平渐近线．
\item 既有铅直又有斜渐近线．
\end{abcd}
\end{problem}


\begin{problem}
设函数$f(x) = \left| x^3 - 1 \right|\phi(x)$，其中$\phi (x)$在$x = 1$处连续，
则$\phi (1) = 0$是$f(x)$在$x = 1$处可导的\pickout{}
\begin{abcd}
\item 充分必要条件．
\item 必要但非充分条件．
\item 充分但非必要条件．
\item 既非充分也非必要条件．
\end{abcd}
\end{problem}


\useproblem{3}{2}{2}%【同试卷三第二(2)题】


\begin{problem}
设矩阵$B = \begin{pmatrix}
  0 & 0 & 1 \\
  0 & 1 & 0 \\
  1 & 0 & 0
\end{pmatrix}$．已知矩阵$A$相似于$B$，则$r(A - 2E)$与$r(A - E)$之和等于\pickout{}
\begin{abcd}
\item $2$．
\item $3$．
\item $4$．
\item $5$．
\end{abcd}
\end{problem}


\begin{problem}
对于任意二事件$A$和$B$，\pickout{}
\begin{abcd}
\item 若$AB \ne \emptyset$，则$A$,$B$一定独立．
\item 若$AB \ne \emptyset$，则$A$,$B$有可能独立．
\item 若$AB = \emptyset$，则$A$,$B$一定独立．
\item 若$AB = \emptyset$，则$A$,$B$一定不独立．
\end{abcd}
\end{problem}


\begin{problem}
设随机变量$X$和$Y$都服从正态分布，且它们不相关，则\pickout{}
\begin{abcd}
\item $X$与$Y$一定独立．
\item $(X,Y)$服从二维正态分布．
\item $X$与$Y$未必独立．
\item $X + Y$服从一维正态分布．
\end{abcd}
\end{problem}


\makepart{}{本题满分8分}%【类似试卷三第三题】

\begin{problem}
设$f(x) = \frac{1}{\sin \pi x} - \frac{1}{\pi x} - \frac{1}{\pi(1 - x)}$，
$x \in\left(0,\frac{1}{2}\right]$．试补充定义$f(0)$，使得$f(x)$在$\left[0,\frac{1}{2}\right]$上连续．
\end{problem}


\makepart{}{本题满分8分}%【同试卷三第四题】

\useproblem{3}{4}{0}


\makepart{}{本题满分8分}%【同试卷三第五题】

\useproblem{3}{5}{0}


\makepart{}{本题满分9分}

\begin{problem}
设$a > 1$，$f(t) = a^t - at$在$(-\infty , +\infty)$内的驻点为$t(a)$．问$a$为何值时，$t(a)$最小？并求出最小值．
\end{problem}


\makepart{}{本题满分9分}

\begin{problem}
设$y = f(x)$是第一象限内连接点$A(0,1)$，$B(1,0)$的一段连续曲线，
$M(x,y)$为该曲线上任意一点，点$C$为$M$在$x$轴上的投影，$O$为坐标原点．
若梯形$OCMA$的面积与曲边三角形$CBM$的面积之和为$\frac{x^3}{6} + \frac{1}{3}$，求$f(x)$的表达式．
\end{problem}


\makepart{}{本题满分8分}

\begin{problem}
设某商品从时刻$0$到时刻$t$的销售量为$x(t) = kt$，$t \in[0,T]$，$k > 0$．
欲在$T$时将数量为$A$的该商品销售完，试求
\begin{enumerate}
\item $t$时刻的商品剩余量，并确定$k$的值；
\item 在时间段$[0,T]$上的平均剩余量．
\end{enumerate}
\end{problem}


\makepart{}{本题满分13分}

\begin{problem}
设有向量组①：$\alpha_1 = (1,0,2)^T$，$\alpha_2 = (1,1,3)^T$，$\alpha_3 = (1, - 1,a + 2)^T$和
向量组②：$\beta_1 = (1,2,a + 3)^T$，$\beta_2 = (2,1,a + 6)^T$，$\beta_3 = (2,1,a + 4)^T$．
试问：当$a$为何值时，向量组①与②等价？当$a$为何值时，向量组①与②不等价？
\end{problem}


\makepart{}{本题满分13分}

\begin{problem}
设矩阵$A = \begin{pmatrix}
  2 & 1 & 1 \\
  1 & 2 & 1 \\
  1 & 1 & a
\end{pmatrix}$可逆，向量$\alpha = \begin{pmatrix}
  1 \\
  b \\
  1
\end{pmatrix}$是矩阵$A^*$的一个特征向量，$\lambda$是$\alpha$对应的特征值，
其中$A^*$是矩阵$A$的伴随矩阵．试求$a$, $b$和$\lambda$的值．
\end{problem}


\makepart{}{本题满分13分}%【同试卷三第十一题】

\useproblem{3}{11}{0}


\makepart{}{本题满分13分}

\begin{problem}
对于任意二事件$A$和$B$，$0 < P(A) < 1,0 < P(B) < 1$，
$$\rho = \frac{P(AB) - P(A)P(B)}{\sqrt{P(A)P(B)P(\bar A)P(\bar B)}}$$
称做事件$A$和$B$的相关系数．
\begin{enumerate}
\item 证明事件$A$和$B$独立的充分必要条件是其相关系数等于零；
\item 利用随机变量相关系数的基本性质，证明$\left| \rho \right| \le 1$．
\end{enumerate}
\end{problem}


\end{document}
